\chapter{Campi Vettoriali Conservativi}
\section{Definizione Campo Vettoriale Conservativo}
Un campo vettoriale $F$ si definisce conservativo se può essere ottenuto esattamente come il gradiente di una funzione.

\vspace{\baselineskip}

Ciò significa che deve esistere una funzione scalare 
\begin{definizione}[]{Condizione per Campo Vettoriale Scalare}
    $$
        U:\R^n \to \R \text{ tale che } F = \nabla U
    $$
    \begin{itemize}
        \item $F$ : \textbf{Campo vettoriale}
        \item $U$ : \textbf{Potenziale} del campo vettoriale $F$
    \end{itemize}
\end{definizione}

\newpage

\section{Teorema dell'indipendenza dal cammino}
Questo teorema stabilisce un legame fondamentale tra la scelta del percorso in un integrale di linea e la natura del campo vettoriale stesso.

\subsection{Ipotesi sul dominio e sul campo}
Affinchp il teorema sia applicabile, è necessario definire l'ambiente geometrico e le proprietà del campo:

\begin{itemize}
    \item Il campo vettoriale $F=F_{1}(x,y)e_{1}+F_{2}(x,y)e_{2}$ \uline{deve essere continuo}.
    \item Il dominio $D$ in cui il compo è definito deve essere \uline{aperto e connesso per archi}.
\end{itemize}

Se le ipotesi elencate sono verificate, il teorema afferma che : \uline{l'indipendenza dell'integrale dal cammino implica che il campo vettoriale $F$ è conservativo}.

\subsection{Applicazione pratica}
Da un punto di vista pratico, questo principio ci dice "come" trova il potenziale.

Sapendo con certezza che il percorso non altera l'integrale, possiamo calcolare la funzione potenziale partendo da un punto fisso $P_{0}(x_{0},y_{0})$ e integrandola fino a un punto generico $P(x,y)$ scegliendo la strada più semplice possibile.

La funzione $U(x,y)$ che ne deriverà sarà il potenziale esatto.

\subsection{Dimostrazione}
I passaggi matematici mirano a costruire la funzione $U$ e a dimostrare che le sue derivate restituiscono le componenti originali del campo.

\begin{enumerate}
    \item \textbf{Passo 1: Definizione del potenziale tramite integrale}
    
    \textbf{Procedura:} Sapendo che l'integrale è indipendente dal cammino, si definisce matematicamente la funzione potenziale $U(x,y,z)$ ponendola uguale all'integrale del campo calcolato da un punto iniziale fisso $P_0$ a un punto generico $P$. Si può scegliere il percorso di integrazione più semplice possibile.
    
    \textbf{Esempio pratico (Calcolo):} Prendiamo il campo vettoriale tridimensionale $\mathbf{F} = 2xy^2z^2 \mathbf{e}_1 + 2yx^2z^2 \mathbf{e}_2 + 2zx^2y^2 \mathbf{e}_3$ e l'origine $(0,0,0)$ come $P_0$. Scegliamo un segmento dritto parametrizzato da $x(t)=xt, y(t)=yt, z(t)=zt$ con $t \in [0,1]$. Sostituendo queste equazioni, il calcolo dell'integrale diventa:
    \[
    \int_{0}^{1} (2x^2y^2z^2 t^5 + 2x^2y^2z^2 t^5 + 2x^2y^2z^2 t^5) dt = 6x^2y^2z^2 \int_{0}^{1} t^5 dt
    \]
    Poiché l'integrale definito di $t^5$ vale $\frac{1}{6}$, l'operazione restituisce esattamente l'equazione del potenziale:
    \[
    U(x,y,z) = x^2y^2z^2
    \]

    \item \textbf{Passo 2: Impostazione dell'incremento orizzontale}
    
    \textbf{Procedura:} Per completare la dimostrazione, bisogna provare algebricamente che derivando la funzione $U$ si ottiene di nuovo il campo di partenza. Si imposta quindi il calcolo dell'incremento orizzontale della funzione: $\Delta U = U(x+\Delta x, y) - U(x,y)$.
    
    \textbf{Esempio pratico (Calcolo):} Grazie all'indipendenza dal cammino, la sottrazione teorica tra l'integrale lungo (fino a $x+\Delta x$) e quello corto (fino a $x$) si riduce algebricamente al calcolo di un unico integrale, ristretto solamente all'ultimo piccolo segmento:
    \[
    \Delta U = \int_{(x,y)}^{(x+\Delta x, y)} (F_1 dx + F_2 dy)
    \]

    \item \textbf{Passo 3: Parametrizzazione dell'integrale sul segmento}
    
    \textbf{Procedura:} Si esegue il calcolo di questo ultimo integrale esprimendo il segmento in forma parametrica in funzione del tempo $t$ (che varia tra $0$ e $1$), inserendo l'incremento $\Delta x$.
    
    \textbf{Esempio pratico (Calcolo):} Inserendo la parametrizzazione $x(t) = x + t\Delta x$ e $y(t) = y$, il termine $y$ risulta costante e la sua derivata è nulla, quindi il differenziale $dy$ scompare dal calcolo. Il differenziale $dx$ diventa $\Delta x dt$. L'espressione si semplifica e l'integrale si calcola esplicitamente come:
    \[
    \Delta U = \Delta x \int_{0}^{1} F_1(x + t\Delta x, y) dt
    \]

    \item \textbf{Passo 4: Limite e verifica della derivata parziale}
    
    \textbf{Procedura:} L'integrale viene risolto tramite il teorema della media. Infine, si imposta il limite del rapporto incrementale per calcolare l'esatta derivata parziale rispetto a $x$.
    
    \textbf{Esempio pratico (Calcolo):} Per il teorema della media, l'integrale appena calcolato equivale a $F_1(x + \bar{t}\Delta x, y)\Delta x$, con un certo valore intermedio $\bar{t} \in [0,1]$. Per ricavare la derivata $\frac{\partial U}{\partial x}$, calcoliamo il rapporto incrementale $\frac{\Delta U}{\Delta x}$: il termine $\Delta x$ presente al numeratore e al denominatore si semplifica perfettamente. Eseguendo il limite per $\Delta x \to 0$:
    \[
    \frac{\partial U}{\partial x} = \lim_{\Delta x \to 0} \frac{F_1(x + \bar{t}\Delta x, y)\Delta x}{\Delta x} = F_1(x,y)
    \]
    Questo dimostra il teorema. In modo analogo si calcolano e si verificano le altre derivate parziali.
\end{enumerate}

\newpage

\section{Cicruitazione di un campo vettoriale conservativo}
\subsection{Definizione}
L'integrale di linea di 2ª specie calcolato lungo una curva orientata chiusa $C$ prende il nome specificato di \textbf{circuitazione} del campo $F$, e per denotarla si utilizza frequentemente il simbolo speciale $\oint_{C} F \cdot t \, d\ell$.

\begin{definizione}{Definizione circuitazione di un campo vettoriale conservativo}
{\Large
    $$
    \oint_{C} F \cdot t \, d \ell = 0
    $$
}
    \begin{itemize}
        \item $\oint_{c}$: È il simbolo dell'integrale di linea di seconda specie calcolato lungo un \textbf{percorso chiuso}.

        L'anellino sul simbolo dell'integrale indica proprio che stiano calcolando una "circuitazione", ovvero stiamo partendo da un punto per ritornare esattamente nello stesso punto dopo aver compiuto un giro.

        \item $F$: È il \textbf{campo vettoriale conservativo}, definito in un certo dominio $D$
        \item $C$: È una \textbf{curva orientata chiusa} arbitraria.
        \item $t$: Rappresenta il \textbf{vettore tangente} alla curva in ogni suo punto.

        Indica la direzione il verso in cui si sta percorrendo fisicamente la curva $C$.
        \item $d\ell$: È l'elemento infinitesimo di lunghezza della curva.
    \end{itemize}
\end{definizione}

\subsection{Come usarlo}
Si usa principalmente in 2 modi:
\begin{itemize}
    \item Se un esercizio ti chiede di calcolare un integrale su un percorso chiuso e sai già che il campo è conservativo, puoi evitare del tutto i calcoli e scrivere che il risultato è 0.
    \item Si può usare come "test di verifica": se calcoli la circuitazione di un campo e ottieni un numero diverso da 0, hai la prova inconfutabile che quel campo non possiede un potenziale globale su quel dominio.
\end{itemize}

\newpage

\subsection{Esempio di calcolo}
Vogliamo dimostrare che 
{\large
$$
    F = (-\frac{y}{x^2+y^2} , \frac{x}{x^2 + y^2})
$$
}

non possiede un potenziale su tutto il piano

\begin{gather}
    C = \text{circonferenza unitaria} \implies x(t) = \cos(t), y(t)= \sin(t) \qquad t \in [0,2\pi] \\
    \int^{2\pi}_{0}[- \frac{\sin(t)}{\cos^2(t)+\sin^2(t)} \cdot \frac{d(\cos(t))}{dt}+ \frac{\cos(t)}{\cos^2(t)+\sin^2(t)} \cdot \frac{d(\sin(t))}{dt}] dt \\
    \int^{2\pi}_{0}[-\sin(t)(-\sin(t)) + \cos(t)(\cos(t))]dt \\
    \int^{2\pi}_{0}(sin^2(t) + \cos^2(t)) dt \\
    \int^{2\pi}_{0} 1dt 0 2\pi \not = 0
\end{gather}

\section{Condizioni di irrotazionalità di un campo vettoriale}
\subsection{Definizione}
Le condizioni di irrotazionalità sono delle identità matematiche che stabiliscono l'uguaglianza tra le derivate parziali "incrociate" delle componenti di un campo vettoriale.

In termini formali, per un campo vettoriali le cui componenti derivate parziali continue, la condizione generale impone che 

\begin{definizione}{Condizioni di irrotazionalità per $\R^n$}
{\large
    $$
    \frac{\partial F_i}{\partial x_j} = \frac{\partial F_j}{\partial x_i}
    $$
}
\begin{itemize}
    \item $F_{i}$ e $F_{j}$: Rappresentano le \textbf{componenti del campo vettoriale $F$}.
    \item $x_{i}$ e $x_{j}$: Rappresentano le \textbf{coordinate spaziali} del dominio in cui è definito il campo.
\end{itemize}
\end{definizione}

La formulazione esatta dipende dalla dimensione dello spazio in cui il campo è definito
\begin{itemize}
    \item In $\R^2$(nel piano) si riduce a un'unica equazione, ovvero: 
    
    \begin{definizione}{Condizioni di irrotazionalità per $\R^2$}
        {\large
            $$
                \frac{\partial F_{1}}{\partial x_{2}} = \frac{\partial F_{2}}{\partial x_{1}}
            $$
        }
    \end{definizione}
    
    \item In $\R^3$(nello spazio) le condizioni da verificare simultaneamente diventano 3:

    \begin{definizione}{Condizioni di irrotazionalità per $\R^3$}
       {\large
       $$
            \frac{\partial F_{1}}{\partial x_{2}} = \frac{\partial F_{2}}{\partial x_{1}} , 
            \frac{\partial F_{1}}{\partial x_{3}} = \frac{\partial F_{3}}{\partial x_{1}},
            \frac{\partial F_{3}}{\partial x_{2}} = \frac{\partial F_{2}}{\partial x_{3}}
        $$
       }
    \end{definizione}
\end{itemize}

\section{chiusura ad esattezza in 1-forma}
\subsection{Definizione}
Nel linguaggio matematico delle forme differenziali, una 1-forma viene indicata con l'espressione:

\begin{definizione}{Espressione chiusura ad esattezza in 1-forma}
    {\LARGE
        $$
            \omega = F_{1}dx_{1} + \dots + F_{n}dx_{n}
        $$
    }
\begin{itemize}
    \item $F_{1}, \dots , F_{n}$: Rappresentano le \textbf{componenti del campo vettoriale} che definiscono la forma.

    Se la forma risulta "esatta" (il che equivale a dire che il campo è conservativo e possiede un potenziale scalare $U$), ciascuna di queste componenti corrisponde all'esatta derivata parziale del potenziale rispetto alla sua variabile $(F_{i}= \frac{\partial U}{\partial x_{i}})$

    \item $dx_{1}, \dots , dx_{n}$: Rappresentano i \textbf{differenziali delle coordinate spaziali}.
\end{itemize}
\end{definizione}

\subsection{Forma Chiusa}
Una forma si dice \textbf{chiusa} se soddisfa le condizioni di irrotazionalità.

Questo significa che deve valere esattamente la formula delle derivate incociate

\subsection{Forma esatta}
Una forma si dice \textbf{esatta} se possiede un potenziale scalare, ovvero se esiste una funzione $U$ tale per cui la forma $\omega$ coincide con il differenziale esatto della funzione $(\omega = dU)$.

In termini pratici, questo vuol dire che ogni componente della forma è la derivata parziale del potenziale rispetto alla corrispondente variabile, ossia $\omega = F_{1}dx_{1} + \dots + F_{n}dx_{n}$.

\uline{\textbf{Una equazione esatta è sempre chiusa}}

\section{Regioni (o domini) semplicemente connesse}
\subsection{Definizione}
\subsubsection{Nel piano $(\R^2)$}
Un insieme aperto si dice semplicemente connesso se è connesso per archi (cioè puoi unire 2 punti qualsiasi con un percorso) e se \uline{ogni curva chiusa appartenente ad esso racchiude una regione piana che è completamente contenuta nel dominio}.

Se ci fosse un buco nel dominio, una curva che gli gira interono racchiuderebbe quel buco violando questa condizione

\subsubsection{Nello spazio $(\R^3)$}
In modo simile, è semplicemente connesso se per ogni curva chiusa nel dominio è possibile trovare \uline{una superficie, completamente contenuta nel dominio, che ammette quella curva come sua frontiera}.

\subsubsection{Generalizzazione in $\R^n$}
Più in generale, un dominio è semplicemente connesso se ogni curva chiusa al suo interno \uline{può essere deformata con continuità fino a contrarsi in un singolo punto} senza mai uscire dal dominio stesso.

\section{Lemma di poincaré}
\subsection{Enunciato}
Sia $F:D \subset \R^n \to \R^n$ \textbf{un campo vettoriale di classe $C^1$} (cioè le componenti $F_{i}$ del campo sono funzioni continue con derivate parziali continue in $D$).

Se il dominio di definizione $D$ è semplicemente connesso, le condizioni di irrotazionalità $(\frac{\partial F_{i}}{\partial x_{j}} = \frac{\partial F_{j}}{\partial x_{i}})$ non sono solo necessarie, ma sono anche sufficienti per l'esistenza di un potenziale del campo in tale dominio.

\section{Formula di green}
\subsection{Definizione}
Per un campo vettoriale piano descritto dalle componenti $F 0 F_{1}(x,y)e_{1} + F_{2}(x,y)e_{2}$, la formula di Green mette in relazione la \textbf{circuitazione del campo vettoriale lungo un cammino chiuso $C$} con un \uline{integrale doppio valutato sulla superficie $\sigma$} racchiusa da tale cammino.

\begin{definizione}{Formula di green}
    {\large
        $$
        \oint_{C} F \cdot t \, d\ell = \iint_{\sigma}(\frac{\partial F_{2}}{\partial x} - \frac{\partial F_{1}}{\partial y}) \, d\sigma
        $$
    }
    \begin{itemize}
        \item $F$: È il \textbf{campo vettoriale piano} in esame, espresso matematicamente come $F = F_{1}(x,y)e_{1} + F_{2}(x,y)e_{2}$.
        \item $F_{1}$ e $F_{2}$: Rappresentano le \textbf{componenti scalari} del campo vettoriale lungo i 2 assi del piano cartesiano.
        \item $x$ e $y$: Sono le \textbf{coordinate spaziali}.
        \item $e_{1}$ e $e_{2}$: Sono i \textbf{versori}.
        \item $C$: Indica la \textbf{curva orientata chiusa} che costituisce la frontiera del dominio.

        Nell'integrale di sinistra $(\oint_{C})$,si calcola la circuitazione proprio lungo questa linea.
        \item $\sigma$: Rappresenta la \textbf{superficie} o il dominio piano semplice che è racchiuso all'interno della curva $C$.
        \item $t$: È il \textbf{vettore tangente} alla curva $C$.
    \end{itemize}
\end{definizione}

\subsection{Ipotesi e Condizioni di Validità}
Affinchè il teorema sia applicabile, devono essere rispettati alcuni requisiti geometrici e analitici:
\begin{itemize}
    \item \textbf{Dominio semplice e frontiera regolare}: La superficie $\sigma$ deve essere un "dominio piano semplice" e la sua frontiera $C$ deve essere formata da curve differenziabili a tratti
    \item \textbf{Continuità}: Le componenti del campo $F_{1}$ e $F_{2}$, cosi' come le loro derivate parziali che compaiono nella formula $(\frac{\partial F_{2}}{\partial x}$ e $\frac{\partial F_{1}}{\partial y})$ devonono essere continue sul dominio
    \item \textbf{La regola dell'orientazione}: Un aspetto critico per il calcolo è il verso di percorrenza della curva chiusa $C$.
    
    Muovendosi lungo il bordo nel senso della tangente, il dominio $\sigma$ deve rimanere sempre alla sinistra dell'osservatore
\end{itemize}
\subsection{Dimostrazione}
\subsubsection*{1. Dimostrazione per la componente $F_1$}
Supponiamo che la regione $\sigma$ sia \textbf{y-semplice}, definita come:
\[ \sigma = \{ (x, y) \in \mathbb{R}^2 \mid a \le x \le b, \, g_1(x) \le y \le g_2(x) \} \]

Vogliamo dimostrare che:
\[ \iint_{\sigma} \left( -\frac{\partial F_1}{\partial y} \right) d\sigma = \oint_{\partial \sigma} F_1 \, dx \]

Calcoliamo l'integrale doppio applicando il Teorema Fondamentale del Calcolo Integrale rispetto alla variabile $y$:
\begin{align*}
\iint_{\sigma} \frac{\partial F_1}{\partial y} \, d\sigma &= \int_a^b \left[ \int_{g_1(x)}^{g_2(x)} \frac{\partial F_1}{\partial y} \, dy \right] dx \\
&= \int_a^b \left[ F_1(x, g_2(x)) - F_1(x, g_1(x)) \right] dx
\end{align*}

Ora calcoliamo l'integrale di linea lungo il bordo $\partial \sigma$ (percorso in senso antiorario):
\begin{itemize}
    \item Sul bordo inferiore $C_1$: $y = g_1(x)$, con $x$ che va da $a$ a $b$. L'integrale è $\int_a^b F_1(x, g_1(x)) \, dx$.
    \item Sul bordo superiore $C_2$: $y = g_2(x)$, con $x$ che va da $b$ a $a$. L'integrale è $\int_b^a F_1(x, g_2(x)) \, dx = -\int_a^b F_1(x, g_2(x)) \, dx$.
\end{itemize}

Sommando i contributi:
\[ \oint_{\partial \sigma} F_1 \, dx = \int_a^b F_1(x, g_1(x)) \, dx - \int_a^b F_1(x, g_2(x)) \, dx \]
Notiamo che questo risultato è esattamente l'opposto dell'integrale doppio calcolato sopra. Segue che:
\[ \oint_{\partial \sigma} F_1 \, dx = -\iint_{\sigma} \frac{\partial F_1}{\partial y} \, d\sigma \]

\subsubsection*{2. Dimostrazione per la componente $F_2$}

In modo del tutto analogo, supponendo che la regione $\sigma$ sia \textbf{x-semplice} (ovvero compresa tra due funzioni di $y$), si dimostra che:
\[ \oint_{\partial \sigma} F_2 \, dy = \iint_{\sigma} \frac{\partial F_2}{\partial x} \, d\sigma \]

\subsubsection*{3. Sintesi Finale}

Sommando le due uguaglianze ottenute, arriviamo alla formulazione finale del teorema:
\[ \oint_{\partial \sigma} (F_1 \, dx + F_2 \, dy) = \iint_{\sigma} \left( \frac{\partial F_2}{\partial x} - \frac{\partial F_1}{\partial y} \right) d\sigma \]
\subsection{Esempio di calcolo}
\begin{gather}
\intertext{Partiamo dalla formula del Teorema di Green, che calcola l'integrale doppio della "densità di rotazione" su una superficie $\sigma$:}
\iint_{\sigma} \left( \frac{\partial F_2}{\partial x} - \frac{\partial F_1}{\partial y} \right) d\sigma \\
%
\intertext{Scegliamo un campo vettoriale specifico utile per calcolare le aree, ad esempio $\mathbf{F}(x, y) = (-y, x)$. Da qui identifichiamo le componenti $F_1 = -y$ e $F_2 = x$. Impostiamo quindi l'uguaglianza tra l'integrale di linea sul bordo $\partial \sigma$ e l'integrale doppio:}
\oint_{\partial \sigma} (-y \, dx + x \, dy) = \iint_{\sigma} \left( \frac{\partial (x)}{\partial x} - \frac{\partial (-y)}{\partial y} \right) d\sigma \\
%
\intertext{Risolviamo le derivate parziali all'interno della parentesi. Derivando $x$ rispetto a $x$ otteniamo $1$, e derivando $-y$ rispetto a $y$ otteniamo $-1$. Sostituendo nella formula avremo:}
\iint_{\sigma} (1 - (-1)) \, d\sigma \\
%
\intertext{Sottraendo i due risultati ($1 + 1 = 2$), l'integrando diventa una costante. Portiamo il $2$ fuori dal segno di integrale:}
2 \iint_{\sigma} d\sigma \\
%
\intertext{Ricordiamo che l'integrale doppio $d\sigma$ nudo e crudo rappresenta l'Area della regione $\sigma$. Se supponiamo che $\sigma$ sia il cerchio di raggio unitario, la sua area è pari a $\pi$. Il calcolo finale risulta quindi:}
2 \cdot \text{Area}(\sigma) = 2\pi
\end{gather}